We compare restricted-Boltzmann-machine (RBM) wavefunctions at hidden densities $\alpha\in\{1,2,4\}$ with a mean-field product state and a Jastrow ansatz on the periodic 1D transverse-field Ising model with $N=8$--$12$ spins, using variational Monte Carlo with plain SGD or stochastic reconfiguration (SR) and exact diagonalisation as the reference. Ten systems, 13 (size, field) tasks, 5 seeds and two starts (symmetric and symmetry-broken initial fields) give 1300 main runs; per seed the lower-energy start is kept and evaluated by exact enumeration. The $\alpha{=}2$ RBM with SR has a lower mean energy error than mean-field and Jastrow with SR in all 13 tasks. Against Jastrow the margin is a factor of 8.55 to 71 for fields $\Gamma$ between 1 and 1.5, where Jastrow is at the best energy that a sampling-free L-BFGS optimisation finds for it, but only 1.47 to 2.37 for $\Gamma\le0.5$, where in three of four tasks the sampling-free Jastrow optimum is below the RBM's Monte Carlo result, so that margin reflects optimisation rather than the ansatz. From the symmetric start alone, Jastrow stays above mean-field at $\Gamma\le0.5$ and all 28 diverged runs (RBM with SR) occur. For the $\alpha{=}2$ RBM, SR has a lower mean error than SGD in all 13 tasks, by a factor of 2.07 at $\Gamma{=}0.25$ and of 68 to 496 for $\Gamma\ge1$. At the critical field the $\alpha{=}4$ RBM with SR is lower than the $\alpha{=}2$ RBM on every seed. Two of five hypotheses were refuted: mean-field energy error does not peak at the critical field, and RBM-SR error there grows with $N$. In the ordered phase the lower-energy states are often symmetry-broken, with fidelity near $0.5$. This is a CPU-scale study of one model with five seeds and reimplemented baselines; it supports no broader claim.
